\documentclass[10pt,a4paper]{amsart}
\usepackage[margin=19mm]{geometry}
\usepackage{amsmath,amssymb,mathtools}
\usepackage[scr=rsfs,cal=euler]{mathalfa}
\usepackage{xfrac}
\usepackage[unicode,hidelinks,bookmarks=false]{hyperref}
\newcommand{\ie}{i.e.\ }
\newcommand{\cf}{cf.\ }
\newcommand{\resp}{respectively\ }
\newcommand{\Subsection}{\subsection}
\providecommand{\nobreakdash}{\nobreak\hbox{-}}
\setlength{\emergencystretch}{2em}
\multlinegap0pt
\usepackage{bm}
\usepackage{bbm}
\usepackage{dsfont}
\usepackage{xspace}
\usepackage{cancel}
\usepackage{mathtools}
\usepackage{amsthm}
\makeatletter
\@ifundefined{thm}{\newtheorem{thm}{Theorem}}{}
\@ifundefined{cor}{\newtheorem{cor}[thm]{Corollary}}{}
\@ifundefined{lem}{\newtheorem{lem}[thm]{Lemma}}{}
\makeatother
\title[Characterizations of zero singular ideal in \'etale groupoid]{Characterizations of zero singular ideal in \'etale groupoid C*-algebras via compressible maps}
\author{Jeremy B. Hume}
\date{Source: arXiv:2509.07262v3 (CC BY 4.0)}
\begin{document}
\maketitle

\noindent\emph{Source and licence.} Characterizations of zero singular ideal in \'etale groupoid C*-algebras via compressible maps. Version arXiv:2509.07262v3 (CC BY 4.0), distributed under the Creative Commons Attribution 4.0 International licence, as recorded in the arXiv licence metadata for this exact version and shown on the abstract page https://arxiv.org/abs/2509.07262v3. Full rights notice: licenses/arxiv-2509.07262-CC-BY-4.0.txt.\par\smallskip

\noindent\textit{Editorial scope.} Counted unit: the Lem whose source label is \texttt{lem:groupidealcharlocwkcont} in the article (source0.tex lines 1048--1058, source label \texttt{lem:groupidealcharlocwkcont}), delivered together with the article's own complete original proof of it (source0.tex lines 1060--1093, one \texttt{proof} environment of 34 lines). No mathematics is written, repaired or re-derived: statement and proof are the article's own text line for line. Required context retained uncounted: thm:locinvar\&compress (lines 653--656); cor:normeqngpd (lines 682--692); eq:norm3 (lines 688--691). Every definition, display and label of those lines is retained unchanged. Omitted from this package and not redistributed: 1--652, 657--681, 692--1047, 1059--1059, 1094--1649. Nothing inside a retained excerpt is omitted or altered: every retained body is delivered exactly as the article's lines are written, so no omission receipt of any kind accompanies this package. Citations: the retained excerpts carry no citation command at all, so no bibliography entry of the article is transcribed and no reference is redistributed. Reference adaptation: none was needed and none was made; every reference inside the delivered unit resolves inside it, so no unresolved reference or citation is produced. Printed slips kept literal: no printed word, symbol, hypothesis or number of the counted statement or of its proof was altered. Layout and class: the article's own class is \texttt{amsart} with packages including geometry, tikz-cd, amssymb, amsrefs, mathrsfs, hyperref, cleveref, bookmark; the wrapper is the lane's pinned \texttt{amsart} editorial wrapper at 10pt on A4, which restates the theorem-like environments the retained text needs and adds the article's own macro definitions that the retained text needs (none), with extra wrapper packages (bm, bbm, dsfont, xspace, cancel, mathtools). The delivered numbering is the unit's own. Assets: no retained span contains a figure environment, a graphics-inclusion command, a PDF-page inclusion, a table or a TikZ picture; no image file is copied into this package, the package declares no asset dependency and no image file is redistributed with it. Licence: version arXiv:2509.07262v3 of this article is distributed under the Creative Commons Attribution 4.0 International licence, as recorded in the arXiv licence metadata for this exact version and shown on the abstract page https://arxiv.org/abs/2509.07262v3. A full rights notice is delivered at licenses/arxiv-2509.07262-CC-BY-4.0.txt. Characters: every retained excerpt is pure ASCII, so the delivered engine, XeLaTeX with its default Latin Modern text font, reports no missing character.
\begin{thm}
    \label{thm:locinvar&compress}
     Let $G$ be an \'etale groupoid, $X\subseteq G^{0}$ a locally invariant set and $\rho$ a pre-C*-norm for $\mathscr{C}_{c}(G)$. For any local groupoid $H$ about $G|_{X}$ with $H^{0} = G^{0}$, the $*$-homomorphism $\eta_{X}:\mathscr{C}_{c}(H)\to \mathscr{C}_{c}(G|_{X})$ is bounded relative to the norms $\rho$ and $\rho(X)$. Consequently, the restriction map $\eta_{X}:\mathscr{C}_{c}(G)\to \mathscr{C}_{c}(G|_{X})$ extends to $\eta_{X}:C^{*}_{\rho}(G)\to C^{*}_{\rho(X)}(G|_{X})$ and is compressible to $C^{*}_{\rho}(H)\subseteq C^{*}_{\rho}(G)$. Furthermore, any approximate unit $(u_{\lambda})$ for $C_{c}(G^{0}\setminus X)$ is an approximate unit for $\ker(C^{*}_{\rho}(H)\to C^{*}_{\rho(X)}(G|_{X}))$.
\end{thm}

\begin{cor}
\label{cor:normeqngpd}
       Let $G$ be an \'etale groupoid, $X\subseteq G^{0}$ a locally invariant set and $\rho$ a pre-C*-norm for $\mathscr{C}_{c}(G)$.
       
       Then, for any approximate unit $(u_{\lambda})$ for $C_{0}(G^{0}\setminus X)$ and any $a\in C^{*}_{\rho}(G)$, we have 
       
       \begin{equation}
       \label{eq:norm3}
           \lim_{\lambda}\|(1-u_{\lambda})a(1-u_{\lambda})\|_{\rho} = \|\eta_{X}(a)\|_{\rho(X)}.
       \end{equation}
\end{cor}

\begin{lem}
\label{lem:groupidealcharlocwkcont}
    Let $G$ be an \'etale groupoid, $x\in G^{0}$ and $\mathcal{X} = \pi^{-1}_{ess}(x)$. Let $H\subseteq G$ be a local groupoid about $G^{x}_{x}$ with $H^{0} = G^{0}$ and $\mathcal{U}$ a neighbourhood basis for $x$ in $G^{0}$. Choose a bisection $U_{g}$ for each $g\in G^{x}_{x}$.
    
    Suppose $S\subseteq \tilde{G}_{ess}^{0}$ is a subset such that the representation $\oplus_{X\in S\cap\pi_{ess}^{-1}(U)}\lambda_{\tilde{G}_{X}}$ is faithful for $C^{*}_{r}(H|_{U}\cdot \pi^{-1}(U))$ , for all $U\in\mathcal{U}$.

    Then, $\ker(C^{*}_{\widehat{r(\mathcal{X})}}(G^{x}_{x})\to C^{*}_{r(\mathcal{X})}(G^{x}_{x})) = 0$ if and only if for every $\epsilon > 0$, finite set $F\subseteq G^{x}_{x}\setminus \{x\}$ and $U\in\mathcal{U}$ with  $U\subseteq \bigcap_{g\in F} r(U_{g})\cap s(U_{g})$, there are vectors $\psi_{i}\in \ell^{2}(G_{x_{i}}^{U}/X_{i})$, where $X_{i}\in S\cap \pi^{-1}(U)$ and $\pi(X_{i}) = x_{i}$, for $i\leq n$ such that

    $$\sum^{n}_{i=1}\langle \psi_{i}, \psi_{i}\rangle = 1\text{ and }|\sum^{n}_{i=1}\langle 1_{U_{g}} * \psi_{i}, \psi_{i}\rangle |\leq \epsilon, \text{ for all }g\in F,$$
    where $1_{U_{g}}$ denotes the characteristic function on $U_{g}$, for $g\in F$. 
\end{lem}

\begin{proof}
 For $f\in \mathscr{C}_{c}(G)$ and $\psi\in \ell^{2}(G_{x}/X)$, we note that $\lambda_{X}(\iota_{ess}(f))(\psi) = f*\psi$, where the right hand side for $\psi = \delta_{hX}$ satisfies $f*\delta_{hX} = \sum_{s(g) = r(h)}f(g)\delta_{ghX}$. Throughout the proof, we make this identification without comment.

    First, we prove the ``only if'' direction. Let $\oplus_{\tau}\pi_{\tau}$ be the GNS representation of $C^{*}_{r(\mathcal{X})}(G^{x}_{x})$, where $\tau$ denotes a state. Then, $\ker(C^{*}_{\widehat{r(\mathcal{X})}}(G^{x}_{x})\to C^{*}_{r(\mathcal{X})}(G^{x}_{x})) = 0$ if and only if $\lambda_{G^{x}_{x}}$ is weakly contained in $\oplus_{\tau}\pi_{\tau}$. Therefore, for every finite set $F\subseteq G^{x}_{x}\setminus \{x\}$ and $\epsilon > 0$, there are states $\tau_{i}$ on $C^{*}_{r(\mathcal{X})}(G^{x}_{x})$ and $a_{i}\in\mathbb{C}[G]$, for $i\leq n$, such that $$| 1 - \sum^{n}_{i=1}\tau_{i}(a^{*}_{i}a_{i})| = |\langle\lambda_{G^{x}_{x}}(\delta_{x})\delta_{x}, \delta_{x}\rangle - \sum^{n}_{i=1}\tau_{i}(a^{*}_{i} \delta^{*}_{x}*a_{i})| = 0$$ and

    \begin{equation}
    \label{eq:wkequivalence1}
        | \sum^{n}_{i = 1}\tau_{i}(a^{*}_{i}\delta^{*}_{g} * a_{i})| = |\langle \lambda_{G^{x}_{x}}(\delta_{g})\delta_{x},\delta_{x}\rangle - \sum^{n}_{i = 1}\tau_{i}(a^{*}_{i}\delta^{*}_{g}* a_{i})|\leq \frac{\epsilon}{2},
    \end{equation}
for all $g\in F$.

Let $\{U_{g}\}_{g\in F}$ be open bisections and $U\in\mathcal{U}$ such that $U\subseteq \bigcap_{g\in F} r(U_{g})\cap s(U_{g})$. Choose $V\in\mathcal{U}$ such that $\overline{V}^{G^{0}}$ is compact and $\overline{V}^{G^{0}}\subseteq U$. Choose $\phi\in C_{c}(U)$ such that $\phi|_{V} = 1$ and define, for $g\in F$, $\phi_{g} = (\phi\circ r)(\phi\circ s)|_{U_{g}}\in C_{c}(U_{g})$.

Let $H\subseteq G$ be the local groupoid as in the hypothesis. Then, $H|_{V} = H_{1}$ and $H|_{U} = H_{2}$ are local groupoids about $G^{x}_{x}$ and $K_{1} = H_{1}\cdot\tilde{G}_{ess}^{0}$, $K_{2} = H_{2}\cdot G^{0}_{ess}$ are local groupoids about $\pi_{ess}^{-1}(x)$. with $H^{0}_{1} = \pi_{ess}^{-1}(V)$, $H^{0}_{ess} = \pi^{-1}_{ess}(U)$. Then, the restriction maps $q_{1}:C^{*}_{r}(K_{1})\to C^{*}_{r(\mathcal{X})}(G^{x}_{x}\cdot \mathcal{X})$ and $q_{2}:C^{*}_{r}(K_{2})\to C^{*}_{r}(G^{x}_{x}\cdot\mathcal{X})$ are *-homomorphisms by Theorem \ref{thm:locinvar&compress}.

Since $C^{*}_{r(\mathcal{X})}(G^{x}_{x})\subseteq C^{*}_{r(\mathcal{X})}(G^{x}_{x}\cdot\mathcal{X})$, we can extend the states $\tau_{i}$, $i\leq n$, to states $\tilde{\tau}_{i}$ on  $C^{*}_{r(\mathcal{X})}(G^{x}_{x}\cdot\mathcal{X})$. Moreover, under this embedding we have $q_{2}(\varphi_{g}) = \delta_{g}$, where $\varphi_{g} = \iota_{ess}(\phi_{g})$, for all $g\in F$. Let $b_{i}\in \mathscr{C}_{c}(K_{1})$ for $i\leq n$ be such that $q_{1}(b_{i}) = a'_{i}$. Then, $b^{*}_{i}\varphi^{*}_{g}b_{i}\in C^{*}_{r}(K_{1})$ and $q_{1}(b^{*}_{i}\varphi^{*}_{g}b_{i}) = a^{*}_{i}\delta^{*}_{g}*a_{i}$ for all $g\in F$ and $i\leq n$. Since $K_{1} = (H|_{V})\cdot \pi_{ess}^{-1}(V)$, by the hypothesis, the representation $\oplus_{X\in S\cap\pi_{ess}^{-1}(V)}\lambda_{\tilde{G}_{X}}$ is faithful for $C^{*}_{r}(K_{1})$, so we can approximate (in the weak* topology) the states $\tilde{\tau}_{i}\circ q_{1}$ by vector states in $\ell^{2}(G^{V}_{\pi(X)}/X)$, $X\in S\cap \pi^{-1}(V)$. Therefore, for every $i\leq n$ there are vectors $\psi'_{ij}\in \ell^{2}(G^{V}_{x_{ij}}/X_{ij})\subseteq \ell^{2}(G^{U}_{x_{ij}}/X_{ij})$ for $j\leq m_{i}$, where $X_{ij}\in S\cap\pi^{-1}(V)\subseteq S\cap \pi^{-1}(U)$ and $\pi(X_{ij}) = x_{ij}$, such that 
\begin{equation}
\label{eq:wkequivalence2}
    \sum_{i,j}\langle \psi'_{ij}, b_{i}^{*}b_{i}\psi'_{ij}\rangle = 1 \text{ and } |\tau_{i}(a^{*}_{i}\delta^{*}_{g}a_{i}) - \sum^{m_{i}}_{j=1}\langle \psi'_{i,j}, (b_{i}^{*}\varphi^{*}_{g}b_{i})\psi'_{i,j}\rangle|\leq \frac{\epsilon}{2n}
\end{equation}
for all $i\leq n$ and $g\in F$.

So, if we re-index $ij$, $i\leq n$, $j\leq m_{i}$ by $k\leq l$, and write $\psi_{k} = b_{i}\psi'_{ij}$, using the inequalities \ref{eq:wkequivalence1} and \ref{eq:wkequivalence2}, we have 

$$\sum^{l}_{k=1}\langle \psi_{k},\psi_{k}\rangle = 1 \text{ and } |\sum^{l}_{k=1}\langle \varphi_{g}*\psi_{k}, \psi_{k}\rangle |\leq \epsilon,$$
for all $g\in F$.

Now, since $\psi_{k}\in \ell^{2}(G^{V}_{x_{k}}/X_{k})$ and $x_{k}\in V$ we have  $\langle \varphi_{g}*\psi_{k}, \psi_{k}\rangle = \langle (\phi_{g}|_{r^{-1}(V)\cap s^{-1}(V)})*\psi_{k}, \psi_{k}\rangle = \langle (1_{U_{g}}|_{r^{-1}(V)\cap s^{-1}(V)})*\psi_{k}, \psi_{k}\rangle = \langle 1_{U_{g}}*\psi_{k}, \psi_{k}\rangle $, for all $k\leq l$. This proves the ``only if'' direction and the ``moreover'' statement.


We now prove the ``if'' direction. For $a\in \mathbb{C}[G^{x}_{x}]$ with $F = \{g\in G:a(g)\neq 0\}\setminus \{x\}$ and bisections $\{V_{g}\}_{g\in F}$, let $V,U\in\mathcal{U}$ such that $V\subseteq \bigcap_{g\in F}r(V_{g})\cap s(V_{g})$, $\overline{U}^{G^{0}}$ is compact, and $\overline{U}^{G^{0}}\subseteq V$, let $\phi\in C_{c}(V)$ be such that $\phi|_{U} = 1$. Set $\phi_{g} = (\phi\circ r)(\phi\circ s)|_{V_{g}}\in C_{c}(V_{g})$, for $g\in F$ and $b^{\phi} = a(x)\phi + \sum_{g\in F}a(g)\phi_{g}$. By the hypothesis, there are vectors $\psi_{k}\in \ell^{2}(G^{U}_{x_{k}}/X_{k})$ for $k\leq l$ such that $\sum^{l}_{k=1} \langle\psi_{k}, \psi_{k}\rangle = 1$ and $|\sum^{l}_{k=1}\langle \psi_{k}, 1_{U_{g}}*\psi_{k}\rangle| \leq \frac{\epsilon}{M}$, for all $g\in F$, where $M = \sum_{g\in F} |a_{g}|$. Since $\langle \psi_{k}, 1_{U_{g}}*\psi_{k}\rangle = \langle \psi_{k}, (1_{U_{g}}|_{r^{-1}(U)\cap s^{-1}(U)})*\psi_{k}\rangle = \langle \psi_{k}, (\phi_{g}|_{r^{-1}(U)\cap s^{-1}(U)})*\psi_{k}\rangle = \langle \psi_{k}, \phi_{g}*\psi_{k}\rangle$ for all $g\in F$ and $k\leq l$ we have, by the triangle ineqality, $|a(x) - \tau(\iota_{ess}(b^{\phi}))|\leq \epsilon$, where $\tau(c) = \sum^{l}_{k=1}\langle \psi_{k}, c*\psi_{k}\rangle$ for $c\in C^{*}_{r}(\tilde{G}_{ess})$. Since $\tau$ is a state on $C^{*}_{r}(\tilde{G}_{ess})$, it follows that $|a(x)| - \epsilon \leq \|\iota_{ess}(b^{\phi})\|$. 

Choose an approximate unit $\phi_{\lambda}\subseteq C_{c}(G^{0}\setminus \{x\})$ such that $\phi_{\lambda}|_{U_{\lambda}} = 0$ for some $U_{\lambda}$ with $\overline{U_{\lambda}}^{G^{0}}$ compact and $\overline{U_{\lambda}}^{G^{0}}\subseteq V$. Then, with $\phi^{\lambda} = (1-u_{\lambda})\phi(1-u_{\lambda})$, we have $b^{\phi^{\lambda}} = (1-u_{\lambda})b^{\phi}(1-u_{\lambda})$. Therefore, we have $|a(x)| - \epsilon \leq \|\iota_{ess}((1-u_{\lambda})b^{\phi}(1-u_{\lambda}))\|$ for all $\lambda$. By Theorem \ref{thm:locinvar&compress}, $(\iota_{ess}(u_{\lambda}) = u_{\lambda}\circ \pi_{ess})$ is an approximate unit for the kernel of some compression of $\eta_{\mathcal{X}}$ to a *-homomorphism. Therefore, by Equation \ref{eq:norm3} in Corollary \ref{cor:normeqngpd}, we have $|a(x)| - \epsilon \leq \|\eta_{\mathcal{X}}(\iota_{ess}(b^{\phi}))\|_{r(\mathcal{X})} = \|a\|_{r(\mathcal{X})}$. Since $\epsilon > 0$ was arbitrary, it follows that $|a(x)|\leq \|a\|_{r(\mathcal{X})}$. Therefore, $a\mapsto a(x)$ defines a $\|\cdot\|_{r(\mathcal{X})}$ bounded linear functional on $\mathbb{C}[G^{x}_{x}]$. As $a\mapsto a(x)$ defines a positive linear functional on $C^{*}(G^{x}_{x})$ and its GNS representation is unitarily equivalent to the left regular representation, it follows that $\ker(C^{*}(G_{x}^{x})\to C_{r(\mathcal{X})}^{*}(G^{x}_{x})))\subseteq \ker(\lambda_{G^{x}_{x}})$ and therefore $\ker(C^{*}_{\widehat{r(\mathcal{X})}}(G^{x}_{x})\to C^{*}_{r(\mathcal{X})}(G^{x}_{x})) = 0$, proving the ``if'' direction and thus the lemma. 
\end{proof}

\end{document}
